% GENERATED FILE. Edit canonical lessons, then run npm run generate:books.
% canonical-source-hash: fnv1a64:9adbc06b470d280d
% canonical-lessons: MW-C229-ujas, MW-C229-ujlo, MW-C229-jhino, MW-C229-dholo, MW-C229-chamakno

\chapter{Light and White}
\label{ch:mw-light-and-white}
\hlchaptermodality{229}

\begin{chapteropening}
\textbf{By the end of this chapter:} \emph{I can describe daylight, bright moonlight and the colour white.}

Complete the last lesson of chapter 229: I can describe daylight, bright moonlight and the colour white.
\end{chapteropening}

\section[\texorpdfstring{\mw{उजास} (ujās)}{उजास (ujās)}]{\mw{उजास} (ujās) — daylight; brightness}

\label{lesson:MW-C229-ujas}



\begin{quote}
Before the new one: say the Marwadi for cool, then the Marwadi for to dry up.
\end{quote}

\subsection*{You'll want to know: \mw{उजास}}
\textbf{\mw{उजास}} — \emph{ujās} — \textquotedblleft{}daylight; brightness\textquotedblright{}.

\begin{quote}
\mw{तड़कै} \mw{उजास} \mw{आवै} \mw{है।}
\end{quote}

\begin{quote}
\emph{taṛkai ujās āvai hai.}
\end{quote}

\begin{quote}
At daybreak the light comes.
\end{quote}

A masculine noun.

\begin{cousinweb}
\textbf{\mw{उजास}} (\emph{ujās}) is inherited from a reconstructed \emph{*ujjhāṣa}, \textquotedblleft{}shining\textquotedblright{}; CDIAL lists Gujarati \emph{ujās}, \textquotedblleft{}brightness, light\textquotedblright{}, under it. CDIAL gives no Marwari form in this entry (CDIAL 1678).
\end{cousinweb}

\begin{grammarlens}[title={What you've built}]
One more word for this chapter.
\end{grammarlens}

\subsection*{Your turn}
\begin{itemize}
\raggedright
  \item \emph{Say it:} \emph{ujās}
  \item \emph{Say it:} \emph{ujās}, once more
  \item \emph{Say it:} \emph{taṛkai ujās āvai hai.} — At daybreak the light comes.
  \item \emph{Recall:} say the Marwadi for cool, then the Marwadi for to dry up, then say \emph{ujās} again
\end{itemize}

\begin{tcolorbox}[breakable,colback=teal!4,colframe=teal!35!black,title={Before you move on}]
What does \mw{उजास} mean? (\textquotedblleft{}Daylight; brightness\textquotedblright{}.) Say it once more.
\end{tcolorbox}
\section[\texorpdfstring{\mw{ऊजळो} (ūjḷo)}{ऊजळो (ūjḷo)}]{\mw{ऊजळो} (ūjḷo) — bright; clear}

\label{lesson:MW-C229-ujlo}



\begin{quote}
Before the new one: say the Marwadi for to dry up, then the Marwadi for daylight.
\end{quote}

\subsection*{You'll want to know: \mw{ऊजळो}}
\textbf{\mw{ऊजळो}} — \emph{ūjḷo} — \textquotedblleft{}bright; clear\textquotedblright{}.

\begin{quote}
\mw{आज} \mw{रात} \mw{चांद} \mw{ऊजळो} \mw{है।}
\end{quote}

\begin{quote}
\emph{āj rāt chā̃d ūjḷo hai.}
\end{quote}

\begin{quote}
Tonight the moon is bright.
\end{quote}

\begin{cousinweb}
\textbf{\mw{ऊजळो}} (\emph{ūjḷo}) is inherited from Sanskrit \emph{ujjvala}, \textquotedblleft{}blazing up, bright\textquotedblright{}; CDIAL lists Old Marwari \emph{ujala}, \textquotedblleft{}bright\textquotedblright{} (CDIAL 1670).
\end{cousinweb}

\begin{grammarlens}[title={What you've built}]
One more word for this chapter.
\end{grammarlens}

\subsection*{Your turn}
\begin{itemize}
\raggedright
  \item \emph{Say it:} \emph{ūjḷo}
  \item \emph{Say it:} \emph{ūjḷo}, once more
  \item \emph{Say it:} \emph{āj rāt chā̃d ūjḷo hai.} — Tonight the moon is bright.
  \item \emph{Recall:} say the Marwadi for to dry up, then the Marwadi for daylight, then say \emph{ūjḷo} again
\end{itemize}

\begin{tcolorbox}[breakable,colback=teal!4,colframe=teal!35!black,title={Before you move on}]
What does \mw{ऊजळो} mean? (\textquotedblleft{}Bright; clear\textquotedblright{}.) Say it once more.
\end{tcolorbox}
\section[\texorpdfstring{\mw{झीणो} (jhīṇo)}{झीणो (jhīṇo)}]{\mw{झीणो} (jhīṇo) — fine (of sand or cloth); thin}

\label{lesson:MW-C229-jhino}



\begin{quote}
Before the new one: say the Marwadi for daylight, then the Marwadi for bright.
\end{quote}

\subsection*{You'll want to know: \mw{झीणो}}
\textbf{\mw{झीणो}} — \emph{jhīṇo} — \textquotedblleft{}fine (of sand or cloth); thin\textquotedblright{}.

\begin{quote}
\mw{थळ} \mw{मांय} \mw{रेत} \mw{झीणी} \mw{है।}
\end{quote}

\begin{quote}
\emph{thaḷ mā̃y ret jhīṇī hai.}
\end{quote}

\begin{quote}
The sand in the desert is fine.
\end{quote}

\begin{cousinweb}
\textbf{\mw{झीणो}} (\emph{jhīṇo}) is inherited from a reconstructed \emph{*jhīna}, \textquotedblleft{}wasted, thin\textquotedblright{}; CDIAL lists Old Marwari \emph{jhīṇa}, \textquotedblleft{}thin\textquotedblright{}, and Gujarati \emph{jhīṇũ}, \textquotedblleft{}thin, fine\textquotedblright{} (CDIAL 5395).
\end{cousinweb}

\begin{grammarlens}[title={What you've built}]
One more word for this chapter.
\end{grammarlens}

\subsection*{Your turn}
\begin{itemize}
\raggedright
  \item \emph{Say it:} \emph{jhīṇo}
  \item \emph{Say it:} \emph{jhīṇo}, once more
  \item \emph{Say it:} \emph{thaḷ mā̃y ret jhīṇī hai.} — The sand in the desert is fine.
  \item \emph{Recall:} say the Marwadi for daylight, then the Marwadi for bright, then say \emph{jhīṇo} again
\end{itemize}

\begin{tcolorbox}[breakable,colback=teal!4,colframe=teal!35!black,title={Before you move on}]
What does \mw{झीणो} mean? (\textquotedblleft{}Fine (of sand or cloth); thin\textquotedblright{}.) Say it once more.
\end{tcolorbox}
\section[\texorpdfstring{\mw{धोळो} (dhoḷo)}{धोळो (dhoḷo)}]{\mw{धोळो} (dhoḷo) — white}

\label{lesson:MW-C229-dholo}



\begin{quote}
Before the new one: say the Marwadi for bright, then the Marwadi for fine.
\end{quote}

\subsection*{You'll want to know: \mw{धोळो}}
\textbf{\mw{धोळो}} — \emph{dhoḷo} — \textquotedblleft{}white\textquotedblright{}.

\begin{quote}
\mw{म्हारो} \mw{कुरतो} \mw{धोळो} \mw{है।}
\end{quote}

\begin{quote}
\emph{mhāro kurto dhoḷo hai.}
\end{quote}

\begin{quote}
My kurta is white.
\end{quote}

\begin{cousinweb}
\textbf{\mw{धोळो}} (\emph{dhoḷo}) is inherited from Sanskrit \emph{dhavala}, \textquotedblleft{}dazzling white\textquotedblright{}; CDIAL lists Hindi \emph{dhaulā}, \textquotedblleft{}white\textquotedblright{}, and a Gujarati cousin. CDIAL gives no Marwari form in this entry (CDIAL 6767).
\end{cousinweb}

\begin{grammarlens}[title={What you've built}]
One more word for this chapter.
\end{grammarlens}

\subsection*{Your turn}
\begin{itemize}
\raggedright
  \item \emph{Say it:} \emph{dhoḷo}
  \item \emph{Say it:} \emph{dhoḷo}, once more
  \item \emph{Say it:} \emph{mhāro kurto dhoḷo hai.} — My kurta is white.
  \item \emph{Recall:} say the Marwadi for bright, then the Marwadi for fine, then say \emph{dhoḷo} again
\end{itemize}

\begin{tcolorbox}[breakable,colback=teal!4,colframe=teal!35!black,title={Before you move on}]
What does \mw{धोळो} mean? (\textquotedblleft{}White\textquotedblright{}.) Say it once more.
\end{tcolorbox}
\section[\texorpdfstring{\mw{चमकणो} (chamakṇo)}{चमकणो (chamakṇo)}]{\mw{चमकणो} (chamakṇo) — to shine; to glitter}

\label{lesson:MW-C229-chamakno}



\begin{quote}
Before the new one: say the Marwadi for fine, then the Marwadi for white.
\end{quote}

\subsection*{You'll want to know: \mw{चमकणो}}
\textbf{\mw{चमकणो}} — \emph{chamakṇo} — \textquotedblleft{}to shine; to glitter\textquotedblright{}.

\begin{quote}
\mw{रात} \mw{तारो} \mw{चमकै} \mw{है।}
\end{quote}

\begin{quote}
\emph{rāt tāro chamkai hai.}
\end{quote}

\begin{quote}
At night a star shines.
\end{quote}

Past (it …): \mw{वो} \mw{तारो} \mw{चमक्यो}, \emph{vo tāro chamkyo}.

\begin{cousinweb}
\textbf{\mw{चमकणो}} (\emph{chamakṇo}) goes back to a reconstructed \emph{*camakka}, \textquotedblleft{}a sudden movement\textquotedblright{}; CDIAL lists Hindi \emph{camaknā} and Gujarati \emph{camakvũ}, \textquotedblleft{}to glitter\textquotedblright{}. CDIAL gives no Marwari form in this entry (CDIAL 4676). \textbf{\mw{चमकीलो}}, shiny, is built on it.
\end{cousinweb}

\begin{grammarlens}[title={What you've built}]
One more word for this chapter.
\end{grammarlens}

\subsection*{Your turn}
\begin{itemize}
\raggedright
  \item \emph{Say it:} \emph{chamakṇo}
  \item \emph{Say it:} \emph{chamakṇo}, once more
  \item \emph{Say it:} \emph{rāt tāro chamkai hai.} — At night a star shines.
  \item \emph{Recall:} say the Marwadi for fine, then the Marwadi for white, then say \emph{chamakṇo} again
\end{itemize}

\begin{tcolorbox}[breakable,colback=teal!4,colframe=teal!35!black,title={Before you move on}]
What does \mw{चमकणो} mean? (\textquotedblleft{}To shine; to glitter\textquotedblright{}.) Say it once more.
\end{tcolorbox}
